\subsection{\texorpdfstring{Properties of mappings \(\Theta_{a}\)}{Properties of mappings \textbackslash Theta\_\{a\}}}

Recall (VAR, R2, pp.~46 and~47, 11.1.3, d)) that if K is a compact subset of C, there exists a fundamental system of compact neighborhoods V of K which are pieces of C (that is, for every \(x \in V\) there exists a chart \((\mathrm{U}, \varphi, \mathbf{C})\) from C at x such that \(\varphi(\mathrm{U} \cap \mathrm{K})\) is an open subset of a closed half-space of C). We then denote by \(\partial V\) the boundary of the piece V equipped with the orientation induced by the orientation of C (VAR, R2, p.~47) and let dz denote the differential of the injection \(\partial V \to C\) .

PROPOSITION 1. --- Let a be an element of A, let U be an open neighborhood of Sp(a), and let \(f \in \mathcal{O}(\mathrm{U};\mathrm{A})\). Let V be a compact neighborhood of Sp(a) contained in U and such that V is a piece of C.

Then \(z \mapsto f(z)(z - a)^{-1}\) is continuous on \(\partial\mathrm{V}\); the differential form \(f(z)(z - a)^{-1}dz\) is integrable on \(\partial\mathrm{V}\), and one has

\[
\Theta_ {a} (f) = \frac {1}{2 i \pi} \int_ {\partial \mathrm{V}} f (z) (z - a) ^ {- 1} d z.
\]

Let \(h\) be a map of \(\mathbf{C}\) into \(\mathbf{C}\) infinitely differentiable, equal to 1 in a neighborhood of \(\mathrm{Sp}(a)\) and with compact support contained in the interior of \(V\) (lemma 1 of I, p.~52). Let \(u\) be a map of \(\mathbf{C}\) to A such that \((h, u)\) is adapted to \(a\) (cf. example 3 of I, p. 53). The associated differential form is \(\omega = du \wedge dz\). It follows that \(f\omega = fdu \wedge dz = d(fu dz)\) since \(f\) is holomorphic. Moreover, \(u(z) = (z - a)^{-1}\) on the boundary of \(V\) . Furthermore, the differential form \(fu dz\) is of class \(C^1\) on U. Therefore,

\[
2 i \pi \Theta_ {a} (f) = \int_ {\mathrm{V}} d (f u d z) = \int_ {\partial \mathrm{V}} f u d z = \int_ {\partial \mathrm{V}} f (z) (z - a) ^ {- 1} d z
\]

by formula (5) and Stokes' formula for the piece V (VAR, R2, p.~47, 11.2.3).

\[
\text { COROLLARY. } - \text {   Let   } a \in \mathrm{A}. \text {   One has   } \Theta_ {a} (1) = 1.
\]

Let \(\mathrm{R} > \varrho(a)\) be a real number. Let V be the closed disk with center 0 and radius R, so that \(\mathrm{Sp}(a) \subset \mathring{\mathrm{V}}\) . This is a piece of C whose boundary \(\partial\mathrm{V}\) is the circle with center 0 and radius R. For \(z \in \mathbf{C} - \mathring{\mathrm{V}}\) , we have the formula \((z - a)^{-1} = z^{-1}(1 - z^{-1}a)^{-1} = \sum_{j=1}^{+\infty} z^{-j} a^{j-1}\) . The series converges uniformly for \(z \in \partial V\) . We therefore have

\[
\Theta_ {a} (1) = \frac {1}{2 i \pi} \sum_ {j = 1} ^ {+ \infty} a ^ {j - 1} \int_ {\partial \mathrm{V}} z ^ {- j} d z = 1
\]

since

\[
\int_ {\partial \mathrm{V}} z ^ {j} d z = 0
\]

for every integer \(j \neq -1\) and

\[
\int_ {\partial \mathrm{V}} z ^ {- 1} d z = 2 i \pi
\]

(VAR, R2, p.~44, 10.4.5, and p.~47, 11.2.1, example).

Lemma 8. --- Let U be an open subset of \(C^{n}\) . Let \(\omega_{1}\) be a continuous differential form of degree n in U with compact support and values in A (resp. let \(\omega_{2}\) be a continuous differential form of degree 2 in C with compact support and values in C). Let us denote by \(\pi_{1}\) and \(\pi_{2}\) the canonical projections of U × C onto U and C. The differential form \(\pi_{1}^{*}\omega_{1} \wedge \pi_{2}^{*}\omega_{2}\) on U × C is continuous with compact support and takes values in A. We have

\[
\int_ {\mathrm{U} \times \mathbf {C}} \pi_ {1} ^ {*} \omega_ {1} \wedge \pi_ {2} ^ {*} \omega_ {2} = \left(\int_ {\mathbf {C}} \omega_ {2}\right) \left(\int_ {\mathrm{U}} \omega_ {1}\right).
\]

Let \(\mu_{n}\) be the Lebesgue measure on \(C^{n}\) and \(\mu_{1}\) the Lebesgue measure on C. There exists \(\psi_{1} \in \mathcal{K}(U;A)\) and \(\psi_{2} \in \mathcal{K}(C)\) such that the vector measure associated with \(\omega_{1}\) is equal to \(\psi_{1} \cdot \mu_{n}\) , and the vector measure associated with \(\omega_{2}\) is equal to \(\psi_{2} \cdot \mu_{1}\) . The vector measure associated with the differential form \(\pi_{1}^{*}\omega_{1} \wedge \pi_{2}^{*}\omega_{2}\) is \((\psi_{1} \otimes \psi_{2}) \cdot \mu_{n} \otimes \mu_{1}\) .

Let \(\ell\) be a continuous linear form on A. By INT, VI, §2, no. 2, def. 2 and the definition of the product measure (INT, III, §4, no. 1, def. 1), it follows

\[
\begin{array}{r l} & {\ell \Big (\int_ {\mathrm{U} \times \mathbf {C}} \pi_ {1} ^ {*} \omega_ {1} \wedge \pi_ {2} ^ {*} \omega_ {2} \Big) = \int_ {\mathrm{U} \times \mathbf {C}} \ell \circ (\psi_ {1} \otimes \psi_ {2}) \mu_ {n} \otimes \mu_ {1}} \\ & {\qquad = \int_ {\mathrm{U} \times \mathbf {C}} \psi_ {2} (z) \ell (\psi_ {1} (x)) d \mu_ {n} (x) d \mu_ {1} (z)} \\ & {\qquad = \Big (\int_ {\mathbf {C}} \psi_ {2} (z) d \mu_ {1} (z) \Big) \Big (\int_ {\mathrm{U}} \ell (\psi_ {1} (x)) d \mu_ {n} (x) \Big)} \\ & {\qquad = \Big (\int_ {\mathbf {C}} \psi_ {2} \mu_ {1} \Big) \ell \Big (\int_ {\mathrm{U}} \psi_ {1} \mu_ {n} \Big),} \end{array}
\]

whence the result (INT VI, loc. cit.).

Lemma 9. --- Let \(\boldsymbol{a} = (a_1, \ldots, a_n) \in \mathbb{A}^n\) . Let \(p \in \mathbf{N}\) and \(\boldsymbol{a}' = (a_{n+1}, \ldots, a_{n+p}) \in \mathbb{A}^p\) . Then we have \(\Theta_{(\boldsymbol{a}, \boldsymbol{a}')} \circ \pi_{n,n+p}^* = \Theta_{\boldsymbol{a}}\) . In particular, we have \(\Theta_{\boldsymbol{a}}(1) = 1\) .

As we have \(\pi_{n,n+p} = \pi_{n,n+1} \circ \cdots \circ \pi_{n+p-1,n+p}\) , it suffices to prove the first assertion when \(p = 1\) , which we now assume. We simply denote \(\pi = \pi_{n,n+1}\) . It then suffices to show that, for every open neighborhood U of \(\mathrm{Sp}^n(\boldsymbol{a})\) , and every function \(f \in \mathcal{O}(\mathrm{U};\mathrm{A})\) , one has \(\Theta_{(\boldsymbol{a},a_{n+1})}(f \circ \pi) = \Theta_{\boldsymbol{a}}(f)\) . Denote by \(g = f \circ \pi\) . Let \(h\) (resp. \(h'\)) be a mapping from \(\mathbf{C}^n\) into \(\mathbf{C}\) (resp. from \(\mathbf{C}\) into \(\mathbf{C}\)), infinitely differentiable, equal to 1 in a neighborhood of \(\mathrm{Sp}^n(\boldsymbol{a})\) (resp. of \(\mathrm{Sp}(\boldsymbol{a}')\)), with compact support contained in U (resp. in \(\mathbf{C}\) ). There exist mappings ( \(u_1, \ldots, u_n\) ) from \(\mathbf{C}^n\) into A, infinitely differentiable, such that the sequence ( \(h, u_1, \ldots, u_n\) ) is adapted to \(\boldsymbol{a}\) (lemma 3 of I, p.~54), and an infinitely differentiable mapping \(u_{n+1}\) of \(\mathbf{C}\) into A such that the pair ( \(h', u_{n+1} \) ) is adapted to \(a_{n+1}\) (loc. cit.)

For \(z \in C^{n}\) and \(z_{n+1} \in C\) , let \(h''(z, z_{n+1}) = h(z)h'(z_{n+1})\) and \(u_{n+1}''(z, z_{n+1}) = h(z)u_{n+1}(z_{n+1})\) . The functions \(h''\) and \(u_{n+1}''\) are infinitely differentiable in \(C^{n+1}\) . The function \(h''\) is equal to 1 in a neighborhood of \(\mathrm{Sp}^{n+1}(a, a_{n+1})\) , and has compact support contained in \(U \times C\) . For every \(\boldsymbol{w} = (z, z_{n+1}) \in \mathbf{C}^{n+1}\) , one has

\[
\begin{array}{c} (z _ {1} - a _ {1}) (u _ {1} \circ \pi) (\boldsymbol {w}) + \dots + (z _ {n} - a _ {n}) (u _ {n} \circ \pi) (\boldsymbol {w}) \\ + (z _ {n + 1} - a _ {n + 1}) u _ {n + 1} ^ {\prime \prime} (\boldsymbol {w}) = 1 - h (\boldsymbol {z}) + h (\boldsymbol {z}) (1 - h ^ {\prime} (z _ {n + 1})) \\ = 1 - h ^ {\prime \prime} (\boldsymbol {w}), \end{array}
\]

which shows that the sequence \((h''_{1}, u_{1} \circ \pi, \ldots, u_{n} \circ \pi, u_{n+1}')\) is adapted to \((\boldsymbol{a}, a_{n+1})\) . Let \(\omega\) be the associated differential form.

The differential form \(du_{1} \wedge dz_{1} \wedge \cdots \wedge du_{n} \wedge dz_{n} \wedge dh\) on \(C^{n}\) has degree \(2n + 1\) , hence is zero. Consequently,

\[
\begin{array}{l} \omega = d (u _ {1} \circ \pi) \wedge d z _ {1} \wedge \dots \wedge d (u _ {n} \circ \pi) \wedge d z _ {n} \wedge d u _ {n + 1} ^ {\prime \prime} \wedge d z _ {n + 1} \\ = (h \circ \pi) d (u _ {1} \circ \pi) \wedge d z _ {1} \wedge \dots \wedge d (u _ {n} \circ \pi) \wedge d z _ {n} \wedge d u _ {n + 1} \wedge d z _ {n + 1}. \end{array}
\]

Since \(g = f \circ \pi\) , formula (5) and lemma 8 imply

\[
\begin{array}{c} \Theta_ {\boldsymbol {a}, \boldsymbol {a} ^ {\prime}} (g) = \frac {(n + 1) !}{(2 i \pi) ^ {n + 1}} \int_ {\mathrm{U} \times \mathbf {C}} g \omega = \frac {(n + 1) !}{(2 i \pi) ^ {n + 1}} \Big (\int_ {\mathbf {C}} d u _ {n + 1} \wedge d z _ {n + 1} \Big) \\ \qquad \qquad \times \Big (\int_ {\mathrm{U}} f h   d u _ {1} \wedge d z _ {1} \wedge \dots \wedge d u _ {n} \wedge d z _ {n} \Big). \end{array}
\]

On the one hand, we have

\[
\int_ {\mathbf {C}} d u _ {n + 1} \wedge d z _ {n + 1} = 2 i \pi \Theta_ {a _ {n + 1}} ^ {\mathbf {C}} (1) = 2 i \pi \cdot 1
\]

by corollary 5. On the other hand, part c) of lemma 4 of I, p.~54 and the fact that the integral of a closed form is zero (VAR, R2, p.~48, 11.2.4) imply

\[
\begin{array}{c} (n + 1) \int_ {\mathrm{U}} f h d u _ {1} \wedge d z _ {1} \wedge \dots \wedge d u _ {n} \wedge d z _ {n} = \\ \int_ {\mathrm{U}} f   d u _ {1} \wedge d z _ {1} \wedge \dots \wedge d u _ {n} \wedge d z _ {n} = \frac {(2 i \pi) ^ {n}}{n !} \Theta_ {\boldsymbol {a}} (f). \end{array}
\]

Thus one obtains

\[
\Theta_ {\boldsymbol {a}, \boldsymbol {a} ^ {\prime}} (g) = \frac {(n + 1) !}{(2 i \pi) ^ {n + 1}} \times \frac {(2 i \pi) ^ {n}}{(n + 1) !} \Theta_ {\boldsymbol {a}} (f) \times 2 i \pi = \Theta_ {\boldsymbol {a}} (f).
\]

Finally, the formula \(\Theta_{\boldsymbol{a}}(1)=1\) follows from the foregoing and from the corollary of Prop. 1.

Lemma 10. --- Let \(\boldsymbol{a} \in \mathrm{A}^n\) . Let \(g\) be a polynomial function on \(\mathbf{C}^n\) with coefficients in A and \(f \in \mathcal{O}(\mathrm{Sp}^n(\boldsymbol{a});\mathrm{A})\) . We have \(\Theta_{\boldsymbol{a}}(gf) = g(\boldsymbol{a})\Theta_{\boldsymbol{a}}(f)\) . In particular, we have \(\Theta_{\boldsymbol{a}}(g) = g(\boldsymbol{a})\) .

By Lemma 9, it suffices to prove the first assertion.

We denote by \(z_{1},\ldots,z_{n}\) the coordinate functions on \(C^{n}\). Since the map \(\Theta_{a}\) is A-linear, it suffices to prove the assertion of the lemma when \(g=z_{1}^{e_{1}}\cdots z_{n}^{e_{n}}\) , where \((e_{1},\ldots,e_{n})\in\mathbf{N}^{n}\). Proceeding by induction on \(e_{1}+\cdots+e_{n}\) , we reduce to the case where there exists an integer i such that \(1\leqslant i\leqslant n\) and \(g=z_{i}\) .

Let U be an open neighborhood of \(\mathrm{Sp}^{n}(\boldsymbol{a})\) . Let \((h,u_{1},\ldots,u_{n})\) be a sequence adapted to \(\boldsymbol{a}\) such that the support of h is contained in U (Lemma 1 of I, p.~52 and Lemma 3 of I, p.~54), and let \(\omega\) be the associated differential form. By lemma 4, a) of I, p.~54, there exists a differential form \(\beta\) such that

\[
(z _ {i} - a _ {i}) \omega = d (h \beta \wedge d z _ {1} \wedge \dots \wedge d z _ {n}).
\]

Consequently, for every function \(f \in \mathcal{O}(\mathrm{U};\mathrm{A})\) , one has

\[
\left(z _ {i} - a _ {i}\right) f \omega = f d \left(h \beta \wedge d z _ {1} \wedge \dots \wedge d z _ {n}\right) = d \left(f h \beta \wedge d z _ {1} \wedge \dots \wedge d z _ {n}\right)
\]

since f is holomorphic, so that \(df \wedge dz_{1} \wedge \cdots \wedge dz_{n} = 0\). Applying Stokes' formula (VAR, R2, p.~48, 11.2.4), we obtain \(\int_{\mathrm{U}}(z_{i} - a_{i})f\omega = 0\) , whence

\[
\Theta_ {\boldsymbol {a}} (z _ {i} f) = \frac {n !}{(2 i \pi) ^ {n}} \int_ {\mathrm{U}} z _ {i} f \omega = \frac {n !}{(2 i \pi) ^ {n}} \int_ {\mathrm{U}} a _ {i} f \omega = a _ {i} \Theta_ {\boldsymbol {a}} (f)
\]

by formula (5). The result follows.

PROPOSITION 2. — Let \(\varrho_{1},\ldots,\varrho_{n}\) be real numbers \(>0\) and let \(\mathrm{U}\subset\mathbf{C}^{n}\) be the polydisc product of the open disks with center 0 and radius \(\varrho_{i}\) . Let

\[
\sum_ {(k _ {1}, \dots , k _ {n}) \in \mathbf {N} ^ {n}} c (k _ {1}, \dots , k _ {n}) \mathrm{X} _ {1} ^ {k _ {1}} \dots \mathrm{X} _ {n} ^ {k _ {n}} \in \mathrm{A} [ [ \mathrm{X} _ {1}, \dots , \mathrm{X} _ {n} ] ]
\]

a formal series with coefficients in A. Suppose that this series converges in U, and denote by f the holomorphic function in U which is its sum.

Let \(\boldsymbol{a}=(a_{1},\ldots,a_{n})\in\mathrm{A}^{n}\) such that \(\varrho(a_{i})<\varrho_{i}\) for \(1\leqslant i\leqslant n\) . Then \(\mathrm{Sp}^{n}(\boldsymbol{a})\subset\mathrm{U}\) , the family \((c(k_{1},\ldots,k_{n})a_{1}^{k_{1}}\cdots a_{n}^{k_{n}})\) of elements of A is absolutely summable, and

\[
\Theta_ {\boldsymbol {a}} (f) = \sum_ {(k _ {1}, \dots , k _ {n}) \in \mathbf {N} ^ {n}} c (k _ {1}, \dots , k _ {n}) a _ {1} ^ {k _ {1}} \dots a _ {n} ^ {k _ {n}}.
\]

For every character \(\chi\) of A and every integer \(i\) such that \(1 \leqslant i \leqslant n\) , one has \(|\chi(a_i)| \leqslant \varrho(a_i) < \varrho_i\) , hence \(\mathrm{Sp}^n(\boldsymbol{a}) \subset \mathrm{U}\) by definition of the simultaneous spectrum. Let \(z_1, \ldots, z_n\) be the restrictions to U of the coordinate functions on \(\mathbf{C}^n\) . Then the family \((c(k_1, \ldots, k_n)z_1^{k_1} \cdots z_n^{k_n})\) is summable in \(\mathcal{O}(\mathrm{U}; \mathrm{A})\) with sum \(f\) . Taking into account Lemma 10 and the continuity of the map \(\Theta_{\boldsymbol{a}}^{\mathrm{U}}\) , the family \((c(k_1, \ldots, k_n)a_1^{k_1} \ldots a_n^{k_n})\) is therefore summable in A with sum \(\Theta_{\boldsymbol{a}}(f)\) . For \(1 \leqslant i \leqslant n\) , let \(\lambda_i\) be a real number such that \(\varrho(a_i) < \lambda_i < \varrho_i\) . There exists \(\mathrm{M}_i < +\infty\) such that \(\|a_i^k\| \leqslant \mathrm{M}_i\lambda_i^k\) for every integer \(k \geqslant 0\) . We then have

\[
\sum_ {(k _ {1}, \dots , k _ {n}) \in \mathbf {N} ^ {n}} \| c (k _ {1}, \dots , k _ {n}) \| \| a _ {1} ^ {k _ {1}} \dots a _ {n} ^ {k _ {n}} \| \leqslant
\]

\[
\mathrm{M} _ {1} \dots \mathrm{M} _ {n} \sum_ {(k _ {1}, \dots , k _ {n}) \in \mathbf {N} ^ {n}} \| c (k _ {1}, \dots , k _ {n}) \| \lambda_ {1} ^ {k _ {1}} \dots \lambda_ {n} ^ {k _ {n}}
\]

which is finite by hypothesis, hence the family \((c(k_{1},\ldots,k_{n})a_{1}^{k_{1}}\cdots a_{n}^{k_{n}})\) is absolutely summable.

COROLLARY. --- Suppose A is nonzero. Let \(a \in C^{n} \subset A^{n}\) . We have \(\operatorname{Sp}_{\mathrm{A}}^{n}(\boldsymbol{a}) = \{\boldsymbol{a}\}\) . For every germ \(f \in \mathcal{O}(\{\boldsymbol{a}\}; \mathrm{A})\) , one has \(\Theta_{\boldsymbol{a}}(f) = f(\boldsymbol{a})\) .

PROPOSITION 3. --- Let B be a commutative unital Banach algebra and \(\varphi\) a continuous unital morphism from A to B. Let \(\boldsymbol{a} = (a_1, \ldots, a_n) \in \mathrm{A}^n\) . Set \(\boldsymbol{b} = (\varphi(a_1), \ldots, \varphi(a_n))\) , so that \(\mathrm{Sp}_{\mathrm{B}}^{n}(\boldsymbol{b}) \subset \mathrm{Sp}_{\mathrm{A}}^{n}(\boldsymbol{a})\) . For every \(f \in \mathcal{O}(\mathrm{Sp}_{\mathrm{A}}^{n}(\boldsymbol{a}); \mathrm{A})\) , one has

\[
\varphi (\Theta_ {\boldsymbol {a}} (f)) = \Theta_ {\boldsymbol {b}} (\varphi_ {*} (f)),
\]

where \(\varphi_{*}(f)\) denotes the germ of \(\varphi \circ f\) in a neighborhood of \(\mathrm{Sp}_{\mathrm{B}}^{n}(\boldsymbol{b})\) .

It suffices to prove that for every open neighborhood U of \(\mathrm{Sp}_{\mathrm{A}}^{n}(\boldsymbol{a})\) and every \(f \in \mathcal{O}(\mathrm{U}; \mathrm{A})\) , one has \(\varphi(\Theta_{\boldsymbol{a}}(f)) = \Theta_{\boldsymbol{b}}(\varphi \circ f)\) , where \(\varphi \circ f \in \mathcal{O}(\mathrm{U}; \mathrm{B})\) . Let \((h, u_{1}, \ldots, u_{n})\) be a sequence adapted to \(\boldsymbol{a}\), where the support of h is contained in U (Lemma 1 of I, p.~52 and Lemma 3 of I, p.~54). Denote by \(\omega\) the associated differential form. For every \(z \in C^{n}\) , one has

\[
\sum_ {j = 1} ^ {n} (z _ {j} - b _ {j}) \varphi (u _ {i} (\boldsymbol {z})) = \varphi \left(\sum_ {j = 1} ^ {n} (z _ {j} - a _ {j}) u _ {j} (\boldsymbol {z})\right) = 1 - h (\boldsymbol {z}),
\]

so that the sequence \((h,\varphi\circ u_{1},\ldots,\varphi\circ u_{n})\) is adapted to b. Let \(\omega'\) be the associated differential form. Denote by \(\mu\) the Lebesgue measure on \(C^{n}\) . Let \(f\in\mathcal{O}(U;A)\) . Write \(\psi\cdot\mu\) for the vector measure associated with the differential form \(f\omega\) . The vector measure associated with the differential form

\[
(\varphi \circ f) \omega^ {\prime} = (\varphi \circ f) d (\varphi \circ u _ {1}) \wedge d z _ {1} \wedge \dots \wedge d (\varphi \circ u _ {n}) \wedge d z _ {n}
\]

is equal to \((\varphi \circ \psi) \cdot \mu\) . Therefore, by formula (5), and INT, VI, §2, no. 2, prop. 2, we have

\[
\Theta_ {\boldsymbol {b}} (\varphi \circ f) = \frac {n !}{(2 i \pi) ^ {n}} \int_ {\mathrm{U}} (\varphi \circ f) \mu = \frac {n !}{(2 i \pi) ^ {n}} \varphi \left(\int_ {\mathrm{U}} \psi \mu\right) = \varphi (\Theta_ {\boldsymbol {a}} (f)),
\]

as was required.

COROLLARY 1. --- Let \(\chi \in \mathrm{X}(\mathrm{A})\) and \(\boldsymbol{a} \in \mathrm{A}^n\) . For every germ \(f \in \mathcal{O}(\mathrm{Sp}^n(\boldsymbol{a}))\) , one has \(\chi(\Theta_{\boldsymbol{a}}(f)) = f(\chi(a_1), \ldots, \chi(a_n))\) .

This is a consequence of Proposition 3, applied to the continuous unital morphism \(\chi\colon\mathbf{A}\to\mathbf{C}\) (th. 1 of I, p.~29), and of the corollary of prop. 2, applied to the Banach algebra C.

Remark. --- Suppose that the algebra A is semisimple. Let \(\boldsymbol{a} = (a_{1}, \ldots, a_{n}) \in \mathrm{A}^{n}\) . By prop. 8 of I, p.~38, the map \(\Theta_{\boldsymbol{a}}^{\mathrm{U}}\) is the unique map \(\varphi\) of \(\mathcal{O}(\mathrm{U})\) to A such that \(\chi(\varphi(f)) = f(\chi(a_{1}), \ldots, \chi(a_{n}))\) for all \(\chi \in \mathrm{X}(\mathrm{A})\) and every function \(f \in \mathcal{O}(\mathrm{U})\) .

COROLLARY 2. --- Let p be an integer ≥ 1. For every family \((f_{1},\ldots,f_{p})\) of elements of \(\mathcal{O}(\mathrm{Sp}^{n}(\boldsymbol{a}))\) , one has

\[
\operatorname{Sp} ^ {p} \left(\left(\Theta_ {\boldsymbol {a}} \left(f _ {1}\right), \dots , \Theta_ {\boldsymbol {a}} \left(f _ {p}\right)\right)\right) = \left(f _ {1}, \dots , f _ {p}\right) \left(\operatorname{Sp} ^ {n} (\boldsymbol {a})\right).
\]

In particular, for every \(f \in \mathcal{O}(\mathrm{Sp}^n(\boldsymbol{a}))\) , one has \(\mathrm{Sp}(\Theta_{\boldsymbol{a}}(f)) = f(\mathrm{Sp}^n(\boldsymbol{a}))\) .

This follows from Cor. 1 and the definition of the simultaneous spectrum.

Example. --- Let A be the complex Banach algebra of functions on the unit circle with absolutely convergent Fourier series (I, p.~19, example 8). Let \(\varphi \in\) A. Let \(f\) be a germ of a holomorphic function in a neighborhood of the set of values of \(\varphi\) . Then \(\psi = \Theta_{\varphi}(f)\) is an absolutely convergent Fourier series which for every \(u \in \mathbf{U}\) satisfies \(\psi(u) = f(\varphi(u))\) (cor. 1, applied to the characters \(\varphi \mapsto \varphi(u)\) ). In other words, the function \(f \circ \varphi\) on the unit circle also has an absolutely convergent Fourier series (“P. Lévy’s theorem”). This result generalizes Wiener’s theorem (I, p.~38, example 4), which concerns the case of the function \(f(z) = 1/z\) on \(\mathbf{C} - \{0\}\) when \(\varphi\) does not vanish.

